% problem-000348 2 稀土串级萃取
\section*{2. 稀土串级萃取}
\textit{下列为分问。}

\subsection*{2-1}
⽬前在稀⼟分离⼯业中⼴泛应⽤国产萃取剂P-507，它的化学名称为(2-⼄基⼰基)膦酸(2-⼄基⼰基)酯，画出其化学结构简式，写出萃取三价稀⼟离⼦的化学反应⽅程式（萃取剂⽤HA表⽰，稀⼟离⼦⽤RE3+表⽰）。

\subsection*{2-2}
设萃取反应平衡常数为 $K _ { \mathrm { e x } }$ ，HA的解离常数为 $K _ { \mathfrak { a } }$ ，配合物在⽔相中的稳定常数 $\cdot \beta _ { n }$ ，配合物和萃取剂在有机相和⽔相间的分配常数分别为 $D _ { \mathrm { m } }$ 和 $P _ { \mathrm { n } }$ ，试推导⽤诸常数表达 $K _ { \mathrm { e x } }$ 的关系式。

\subsection*{2-3}
图2.1是⽤酸性磷酸酯萃取镨离⼦(Pr3+)得到的配合物结构的⼀部分。

（图：）
图 2.1

\subsection*{2-3}
-1 两图的 R 不同，若 R 为 CH_{3}和 ${ \mathrm { C } } _ { 2 } { \mathrm { H } } _ { 5 } ,$ ，写出(a)和(b)的化学式。

\subsection*{2-3}
-2 若烷酯基上碳链增⼤，萃合物将取(a)结构还是(b)结构？请简述理由。

\subsection*{2-3}
-3 在稀⼟分离⼯业中，⼀般应保持萃取剂与稀⼟的⽐例在 6:1 以上以避免乳化现象，为什么？

\subsection*{2-4}
⽯油⼯业的副产品环烷酸是⼯业上分离稀⼟的另⼀重要萃取剂，其通式为RCOOH（R为环烷基）。

\subsection*{2-4}
-1 图 2.2 ⽰出镧(La)和钇(Y)的选择性关系。图中 $\beta _ { \mathrm { L a / Y } }$ 为分离系数，即镧的萃取分配⽐和钇的萃取分配⽐的

⽐值，�为平衡⽔相中镧和钇的物质的量之⽐。请回答，当�值增⼤或减⼩时何种元素容易被萃取。（注：分配⽐ $D =$ $c _ { ( 0 ) } / c _ { ( \mathrm { W } ) }$ ，�为被萃取物总浓度）2-4-2 ⾼纯氧化钇是彩⾊电视和三基⾊荧光灯的红⾊荧光粉基质材料，过去⽤两步萃取技术获得，后徐光宪等⽤环烷酸体系经⼀步萃取分离就可获得99.99%的⾼纯氧化钇，创⽴了国际领先⽔平的新⼯艺。该⼯艺需在煤油中添加 $1 5 { \sim } 2 0 \%$ (V/V)的混合醇，如正庚醇和正癸醇，请简要说明混合醇对改善萃取剂结构性能和萃取作业所起的作⽤。

（图：）
图 2.2
